\documentclass{beamer}
\usetheme{Madrid}
\usecolortheme{default}
\usepackage{amsmath, amssymb, graphicx, hyperref, array, booktabs}
\usepackage[utf8]{inputenc}
\usepackage{color}

% Define colors
\definecolor{cfablue}{RGB}{0,51,102}
\definecolor{cfagold}{RGB}{204,153,0}
\setbeamercolor{title}{fg=cfablue}
\setbeamercolor{frametitle}{fg=cfablue}
\setbeamercolor{structure}{fg=cfablue}

% Title information
\title{The Relationship Between Legality and Ethics in the Finance Profession}
\subtitle{CFA Institute Standards of Practice Handbook Analysis}
\author{Scholar Practioner Aspect Candidate}
\date{\today}

\begin{document}
	
	% ============================================================
	% SLIDE 1: TITLE
	% ============================================================
	\begin{frame}
		\titlepage
	\end{frame}
	
	% ============================================================
	% SLIDE 2: AGENDA
	% ============================================================
	\begin{frame}{Agenda}
		\begin{enumerate}
			\item \textbf{Core Question}: Does Illegal Always Mean Unethical?
			\item \textbf{CFA Framework}: The "More Strict" Standard
			\item \textbf{Young's Argument}: Valid Points or Justification?
			\item \textbf{Challenging Scenarios} (10 Hard Questions)
			\item \textbf{Key Takeaways} for DBA Scholars
		\end{enumerate}
	\end{frame}
	
	% ============================================================
	% SLIDE 3: THE CENTRAL TENSION
	% ============================================================
	\begin{frame}{The Central Tension}
		\begin{block}{The Question}
			Does illegal always mean unethical? More importantly, does legal always mean ethical?
		\end{block}
		
		\begin{columns}
			\column{0.5\textwidth}
			\textbf{Common Assumption}
			\begin{itemize}
				\item Illegal = Unethical
				\item Legal = Ethical
				\item Compliance = Sufficient
			\end{itemize}
			
			\column{0.5\textwidth}
			\textbf{CFA Position}
			\begin{itemize}
				\item Not necessarily
				\item Legal is baseline, not ceiling
				\item \textcolor{cfablue}{"More strict" standard applies}
			\end{itemize}
		\end{columns}
		
		\vspace{0.3cm}
		\footnotesize \textit{Source: CFA Institute Standards of Practice Handbook (2014, p. 14)}
	\end{frame}
	
	% ============================================================
	% SLIDE 4: STANDARD I(A) - KNOWLEDGE OF THE LAW
	% ============================================================
	\begin{frame}{Standard I(A): Knowledge of the Law}
		\begin{block}{The Standard}
			"Members and Candidates must understand and comply with all applicable laws, rules, and regulations... In the event of conflict, Members and Candidates must comply with the more strict law, rule, or regulation."
		\end{block}
		
		\begin{itemize}
			\item \textbf{Key Principle}: Legal compliance is the \textcolor{cfablue}{minimum}
			\item \textbf{Global Application}: Different jurisdictions, different rules
			\item \textbf{The "More Strict" Test}: Compare law vs. Code \& Standards
		\end{itemize}
		
		\vspace{0.3cm}
		\begin{center}
			\fcolorbox{cfablue}{white}{\parbox{0.8\textwidth}{\centering \textbf{When in conflict, follow the HIGHER standard}}}
		\end{center}
	\end{frame}
	
	% ============================================================
	% SLIDE 5: GLOBAL APPLICATION - EXHIBIT 1
	% ============================================================
	\begin{frame}{Global Application of the Code and Standards}
		\begin{table}[]
			\centering
			\small
			\begin{tabular}{|p{3.2cm}|p{3.5cm}|}
				\hline
				\textbf{Applicable Law} & \textbf{Duty} \\
				\hline
				No securities laws & Adhere to Code \& Standards \\
				Less strict than CFA & Adhere to Code \& Standards \\
				More strict than CFA & Adhere to local law \\
				\hline
			\end{tabular}
		\end{table}
		
		\vspace{0.3cm}
		\textbf{Implication}: Legal requirements vary; ethical obligations are universal.
		
		\vspace{0.2cm}
		\footnotesize \textit{Source: CFA Institute (2014, pp. 16-18)}
	\end{frame}
	
	% ============================================================
	% SLIDE 6: YOUNG'S ARGUMENT - OVERVIEW
	% ============================================================
	\begin{frame}{Young's Argument (2019)}
		\begin{block}{Core Thesis}
			"Putting the law in its place" - Challenging the assumption that illegal implies unethical.
		\end{block}
		
		\textbf{Key Points:}
		\begin{enumerate}
			\item Laws can be \textcolor{cfablue}{unjust} or \textcolor{cfablue}{outdated}
			\item Regulatory bodies can be flawed ("patronage haven")
			\item Civil disobedience is ethically defensible
			\item Legal is not Moral; Illegal is not Immoral
		\end{enumerate}
		
		\vspace{0.3cm}
		\textbf{Critical Question}: Is Young exposing flawed ethical reasoning or justifying misconduct?
	\end{frame}
	
	% ============================================================
	% SLIDE 7: HARD QUESTION 1 - TRAVEL FUNDING
	% ============================================================
	\begin{frame}{Hard Question 1: Travel Funding and Independence}
		\begin{block}{The Scenario (CFA Handbook, p. 33)}
			Steven Taylor, a mining analyst, is invited by Precision Metals to tour mining facilities. The company arranges chartered flights and modest accommodations. Taylor allows Precision Metals to pay all expenses.
		\end{block}
		
		\vspace{0.2cm}
		\textbf{Questions:}
		\begin{enumerate}
			\item Is accepting paid travel a violation of Standard I(B)?
			\item If the travel is legal, is it ethical?
			\item What if the accommodations were luxury instead of modest?
			\item Does Young's framework justify Taylor's acceptance?
		\end{enumerate}
		
		\vspace{0.2cm}
		\textcolor{cfablue}{\textbf{CFA Answer}}: Not necessarily a violation IF trip was strictly business and accommodations were modest. However, the \textbf{appearance} of impropriety must be considered.
	\end{frame}
	
	% ============================================================
	% SLIDE 8: HARD QUESTION 2 - GIFTS FROM CLIENTS
	% ============================================================
	\begin{frame}{Hard Question 2: Gifts and Entertainment from Clients}
		\begin{block}{The Scenario (CFA Handbook, pp. 35-36)}
			Theresa Green manages Ian Knowlden's portfolio, achieving superior returns. Knowlden offers Green two Wimbledon tickets and use of his London flat for a week. Green discloses this to her supervisor.
		\end{block}
		
		\vspace{0.2cm}
		\textbf{Questions:}
		\begin{enumerate}
			\item Is accepting this gift a violation?
			\item Does disclosure make it ethical?
			\item What if Green did NOT disclose the gift?
			\item How should the supervisor monitor for favoritism?
		\end{enumerate}
		
		\vspace{0.2cm}
		\textcolor{cfablue}{\textbf{CFA Answer}}: Compliance IF disclosed. BUT supervisor must monitor for pattern of favoritism. \textbf{Transparency is key}.
	\end{frame}
	
	% ============================================================
	% SLIDE 9: HARD QUESTION 3 - EXTERNAL MANAGER TRAVEL
	% ============================================================
	\begin{frame}{Hard Question 3: External Manager Travel Expenses}
		\begin{block}{The Scenario (CFA Handbook, pp. 36-37)}
			Tom Wayne, pension fund manager, recommends Penguin Advisors after attending an "investment fact-finding trip to Asia" partially sponsored by Penguin. Penguin's president was also on the trip.
		\end{block}
		
		\vspace{0.2cm}
		\textbf{Questions:}
		\begin{enumerate}
			\item Was Wayne's recommendation tainted?
			\item What if the trip was educational and the decision was independent?
			\item Does the \textbf{public perception} matter?
			\item What precautions should Wayne have taken?
		\end{enumerate}
		
		\vspace{0.2cm}
		\textcolor{cfablue}{\textbf{CFA Answer}}: Even if no actual violation, the \textbf{perception} of impropriety is problematic. Best practice: employer should pay; contact should be limited; trip should be public record.
	\end{frame}
	
	% ============================================================
	% SLIDE 10: HARD QUESTION 4 - ISSUER-PAID RESEARCH
	% ============================================================
	\begin{frame}{Hard Question 4: Issuer-Paid Research and Compensation}
		\begin{block}{The Scenario (CFA Handbook, p. 37)}
			Javier Herrero is hired by an investor-relations firm to write a research report. The firm will pay a flat fee PLUS a bonus if new investors buy stock as a result of his report.
		\end{block}
		
		\vspace{0.2cm}
		\textbf{Questions:}
		\begin{enumerate}
			\item Is the bonus structure ethical?
			\item Does legality matter if the compensation creates a conflict?
			\item What compensation structure WOULD be ethical?
			\item Does Young's framework justify this arrangement?
		\end{enumerate}
		
		\vspace{0.2cm}
		\textcolor{cfablue}{\textbf{CFA Answer}}: Violates Standard I(B). Bonus creates incentive to draft positive reports regardless of facts. Only flat fee NOT tied to conclusions is acceptable.
	\end{frame}
	
	% ============================================================
	% SLIDE 11: HARD QUESTION 5 - PAY-TO-PLAY SCANDAL
	% ============================================================
	\begin{frame}{Hard Question 5: The "Pay-to-Play" Scandal}
		\begin{block}{The Scenario (CFA Handbook, p. 39)}
			Adrian Mandel, CFA, makes political contributions to Ernesto Gomez, chair of a pension fund investment board. Mandel also entertains Gomez lavishly at top restaurants. All expenses are reimbursed as "marketing expenses."
		\end{block}
		
		\vspace{0.2cm}
		\textbf{Questions:}
		\begin{enumerate}
			\item Are political contributions unethical if legal?
			\item Does disclosure to senior management make it acceptable?
			\item How does this compromise Gomez's objectivity?
			\item What does Young's framework suggest about this case?
		\end{enumerate}
		
		\vspace{0.2cm}
		\textcolor{cfablue}{\textbf{CFA Answer}}: Violates Standard I(B). Even if structured to be technically legal, the conduct is unethical because it compromises objectivity and undermines fair selection processes.
	\end{frame}
	
	% ============================================================
	% SLIDE 12: HARD QUESTION 6 - MATERIAL NONPUBLIC INFORMATION
	% ============================================================
	\begin{frame}{Hard Question 6: Material Nonpublic Information}
		\begin{block}{The Scenario (CFA Handbook, p. 68)}
			Samuel Peter attends a conference call with investment bankers and a company CEO. Salespeople and portfolio managers walk in and out, overhearing that earnings projections have dropped. Before the call ends, they trade on this information.
		\end{block}
		
		\vspace{0.2cm}
		\textbf{Questions:}
		\begin{enumerate}
			\item Did Peter violate Standard II(A)?
			\item What if the information was accidentally overheard?
			\item Does "I didn't know" excuse the conduct?
			\item What firewalls should have been in place?
		\end{enumerate}
		
		\vspace{0.2cm}
		\textcolor{cfablue}{\textbf{CFA Answer}}: Peter violated Standard II(A) by failing to prevent transfer of nonpublic information. The salespeople and portfolio managers also violated by trading on inside information.
	\end{frame}
	
	% ============================================================
	% SLIDE 13: HARD QUESTION 7 - MARKET MANIPULATION
	% ============================================================
	\begin{frame}{Hard Question 7: Market Manipulation}
		\begin{block}{The Scenario (CFA Handbook, pp. 79-80)}
			Bill Mandeville uses mostly positive scenarios as model inputs to achieve the best possible rating for structured products. The outputs support the idea that products have minimal downside risk. Compensation is tied to ratings and sales.
		\end{block}
		
		\vspace{0.2cm}
		\textbf{Questions:}
		\begin{enumerate}
			\item Is manipulating model inputs a violation?
			\item What if the outputs were technically accurate?
			\item Does the intent matter if the market eventually crashes?
			\item How does Young's framework apply to "legal" manipulation?
		\end{enumerate}
		
		\vspace{0.2cm}
		\textcolor{cfablue}{\textbf{CFA Answer}}: Violates Standard II(B). Manipulating inputs to minimize risk misleads market participants and damages trust in capital markets.
	\end{frame}
	
	% ============================================================
	% SLIDE 14: HARD QUESTION 8 - SUITABILITY
	% ============================================================
	\begin{frame}{Hard Question 8: Suitability and Risk Profile}
		\begin{block}{The Scenario (CFA Handbook, p. 108)}
			Caleb Smith recommends investing 20\% of two clients' portfolios in high-tech, high-risk equities. Larry Robertson (60, high net worth) and Gabriel Lanai (40, saving for children's education) have very different risk tolerances.
		\end{block}
		
		\vspace{0.2cm}
		\textbf{Questions:}
		\begin{enumerate}
			\item Is the same recommendation suitable for both clients?
			\item What if both clients \textbf{asked} for this investment?
			\item What is Smith's duty to educate clients?
			\item Does legality matter if the investment is unsuitable?
		\end{enumerate}
		
		\vspace{0.2cm}
		\textcolor{cfablue}{\textbf{CFA Answer}}: Violates Standard III(C) for Lanai. Suitable for Robertson (aggressive profile) but not for Lanai (needs steady returns). Suitability depends on client circumstances.
	\end{frame}
	
	% ============================================================
	% SLIDE 15: HARD QUESTION 9 - PERFORMANCE PRESENTATION
	% ============================================================
	\begin{frame}{Hard Question 9: Performance Presentation}
		\begin{block}{The Scenario (CFA Handbook, pp. 114-115)}
			Kyle Taylor states in marketing materials: "You can expect steady 25\% annual compound growth of the value of your investments." The fund achieved 25\% for ONE year but averages 5\% over five years.
		\end{block}
		
		\vspace{0.2cm}
		\textbf{Questions:}
		\begin{enumerate}
			\item Is the statement a violation?
			\item What if Taylor said "past performance" instead of "can expect"?
			\item Does the statement constitute a guarantee?
			\item What disclosures are required?
		\end{enumerate}
		
		\vspace{0.2cm}
		\textcolor{cfablue}{\textbf{CFA Answer}}: Violates Standard III(D). 25\% growth occurred only in one year, and the statement guarantees future performance - a violation of Standard I(C) as well.
	\end{frame}
	
	% ============================================================
	% SLIDE 16: HARD QUESTION 10 - CONFIDENTIALITY
	% ============================================================
	\begin{frame}{Hard Question 10: Preservation of Confidentiality}
		\begin{block}{The Scenario (CFA Handbook, p. 121)}
			Sarah Connor receives internal reports about City Medical Center's renovation plans. A local businessman asks for this information to decide where to make a substantial contribution. Connor refuses to divulge the information.
		\end{block}
		
		\vspace{0.2cm}
		\textbf{Questions:}
		\begin{enumerate}
			\item Did Connor act correctly?
			\item What if the information would benefit the hospital?
			\item What if the businessman is a major potential donor?
			\item When is disclosure of confidential information permitted?
		\end{enumerate}
		
		\vspace{0.2cm}
		\textcolor{cfablue}{\textbf{CFA Answer}}: Connor correctly refused. Standard III(E) requires confidentiality unless (1) information concerns illegal activities, (2) disclosure required by law, or (3) client permits disclosure.
	\end{frame}
	
	% ============================================================
	% SLIDE 17: HARD QUESTION 11 - REFERRAL FEES
	% ============================================================
	\begin{frame}{Hard Question 11: Referral Fees}
		\begin{block}{The Scenario (CFA Handbook, pp. 203-204)}
			White accepts a new account referred by Brady Securities. Brady receives research reports and directs stock commission business from referral accounts to Brady. White does not disclose this arrangement to the new client.
		\end{block}
		
		\vspace{0.2cm}
		\textbf{Questions:}
		\begin{enumerate}
			\item Is White required to disclose the referral arrangement?
			\item What if the referral was from a trusted friend?
			\item Does the client's satisfaction with the service matter?
			\item What if the fee is paid in services rather than cash?
		\end{enumerate}
		
		\vspace{0.2cm}
		\textcolor{cfablue}{\textbf{CFA Answer}}: Violates Standard VI(C). Disclosure required BEFORE entering into any formal agreement for services. Consideration includes cash, soft dollars, or in-kind.
	\end{frame}
	
	% ============================================================
	% SLIDE 18: HARD QUESTION 12 - DILIGENCE AND REASONABLE BASIS
	% ============================================================
	\begin{frame}{Hard Question 12: Diligence and Reasonable Basis}
		\begin{block}{The Scenario (CFA Handbook, pp. 160-161)}
			Helen Hawke's firm anticipates a tax loophole closing. To capitalize, she estimates IPO prices based on relative company size rather than thorough research, believing she can "justify pricing later."
		\end{block}
		
		\vspace{0.2cm}
		\textbf{Questions:}
		\begin{enumerate}
			\item Is estimating price by size a violation?
			\item What if time pressure makes full research impossible?
			\item Should Hawke have refused the business?
			\item Does Young's framework justify taking on "too much" work?
		\end{enumerate}
		
		\vspace{0.2cm}
		\textcolor{cfablue}{\textbf{CFA Answer}}: Violates Standard V(A). Sarkozi should have taken on only work it could adequately handle. Bypassing research violates due diligence requirements.
	\end{frame}
	
	% ============================================================
	% SLIDE 19: HARD QUESTION 13 - COMMUNICATION WITH CLIENTS
	% ============================================================
	\begin{frame}{Hard Question 13: Communication with Clients}
		\begin{block}{The Scenario (CFA Handbook, p. 173)}
			Sarah Williamson plans to publish an investment newsletter with only weekly "buy" and "sell" recommendations, omitting details of the valuation models and portfolio structuring scheme.
		\end{block}
		
		\vspace{0.2cm}
		\textbf{Questions:}
		\begin{enumerate}
			\item Is omitting model details a violation?
			\item What if the model is proprietary?
			\item Can Williamson provide the information upon request?
			\item What is the minimum required disclosure?
		\end{enumerate}
		
		\vspace{0.2cm}
		\textcolor{cfablue}{\textbf{CFA Answer}}: Violates Standard V(B). Need not describe every detail, but must inform clients of the basic process and logic. Clients must understand limitations and inherent risks.
	\end{frame}
	
	% ============================================================
	% SLIDE 20: HARD QUESTION 14 - LOYALTY TO EMPLOYER
	% ============================================================
	\begin{frame}{Hard Question 14: Loyalty to Employer}
		\begin{block}{The Scenario (CFA Handbook, p. 130)}
			Samuel Magee, frustrated with working environment, is offered a position with a competitor. Before resigning, he asks four big accounts to leave Trust Assets and open accounts with Fiduciary Management.
		\end{block}
		
		\vspace{0.2cm}
		\textbf{Questions:}
		\begin{enumerate}
			\item Did Magee violate Standard IV(A)?
			\item What if he only contacted clients who were "his" accounts?
			\item What if he used public information to find them?
			\item When is soliciting former clients permissible?
		\end{enumerate}
		
		\vspace{0.2cm}
		\textcolor{cfablue}{\textbf{CFA Answer}}: Violates Standard IV(A). Solicitation of current clients and prospective clients before leaving violates duty of loyalty. Must wait until employment ends.
	\end{frame}
	
	% ============================================================
	% SLIDE 21: HARD QUESTION 15 - SUPERVISOR RESPONSIBILITY
	% ============================================================
	\begin{frame}{Hard Question 15: Supervisor Responsibility}
		\begin{block}{The Scenario (CFA Handbook, pp. 148-149)}
			Jane Mattock, head of research, orally advises executives of a proposed recommendation change before the report is published. An executive immediately sells stock from personal and discretionary accounts. Others inform institutional clients before dissemination.
		\end{block}
		
		\vspace{0.2cm}
		\textbf{Questions:}
		\begin{enumerate}
			\item Did Mattock violate Standard IV(C)?
			\item What if she told executives to keep the information confidential?
			\item What compliance procedures should have been in place?
			\item Does delegating responsibility excuse the supervisor?
		\end{enumerate}
		
		\vspace{0.2cm}
		\textcolor{cfablue}{\textbf{CFA Answer}}: Violates Standard IV(C). Failed to prevent dissemination of or trading on unpublished recommendation changes. Must ensure adequate compliance procedures exist.
	\end{frame}
	
	% ============================================================
	% SLIDE 22: HARD QUESTION 16 - ADDITIONAL COMPENSATION
	% ============================================================
	\begin{frame}{Hard Question 16: Additional Compensation}
		\begin{block}{The Scenario (CFA Handbook, p. 140)}
			Geoff Whitman manages Carol Cochran's account. Cochran proposes: "Any year the portfolio achieves at least 15\% return, you and your wife can fly to Monaco at my expense and use my condominium." Whitman accepts without informing his employer.
		\end{block}
		
		\vspace{0.2cm}
		\textbf{Questions:}
		\begin{enumerate}
			\item Is Whitman required to inform his employer?
			\item What if the compensation is "non-monetary"?
			\item Does the conditional nature matter?
			\item What if Whitman's employer permits such arrangements?
		\end{enumerate}
		
		\vspace{0.2cm}
		\textcolor{cfablue}{\textbf{CFA Answer}}: Violates Standard IV(B). Must obtain written consent from all parties. Non-monetary compensation creates conflict and requires disclosure.
	\end{frame}
	
	% ============================================================
	% SLIDE 23: HARD QUESTION 17 - SOCIAL MEDIA
	% ============================================================
	\begin{frame}{Hard Question 17: Social Media and Ethics}
		\begin{block}{The Scenario (CFA Handbook, p. 102)}
			Mary Burdette tweets a "buy" recommendation for Sun Drive Auto stock to FIM Twitter followers before distributing her complete research report to all clients.
		\end{block}
		
		\vspace{0.2cm}
		\textbf{Questions:}
		\begin{enumerate}
			\item Does this violate fair dealing?
			\item What if the tweet was a "summary" rather than a recommendation?
			\item What if all clients follow FIM on Twitter?
			\item How should firms manage social media communications?
		\end{enumerate}
		
		\vspace{0.2cm}
		\textcolor{cfablue}{\textbf{CFA Answer}}: Violates Standard III(B) and IV(C). Investment recommendations must be distributed fairly. Social media does not remove the duty to treat all clients fairly.
	\end{frame}
	
	% ============================================================
	% SLIDE 24: HARD QUESTION 18 - RECORD RETENTION
	% ============================================================
	\begin{frame}{Hard Question 18: Record Retention}
		\begin{block}{The Scenario (CFA Handbook, p. 183)}
			Martin Blank develops an analytical model while employed at GPIM. He takes copies of all supporting records to a new employer without GPIM's permission.
		\end{block}
		
		\vspace{0.2cm}
		\textbf{Questions:}
		\begin{enumerate}
			\item Does Blank violate Standard V(C)?
			\item What if Blank developed the model on his own time?
			\item What if he only takes the "ideas" but not the actual records?
			\item How should Blank use the model in the future?
		\end{enumerate}
		
		\vspace{0.2cm}
		\textcolor{cfablue}{\textbf{CFA Answer}}: Violates Standard V(C). Records created as part of professional activity are the property of the firm. Blank must re-create records at the new firm from public sources.
	\end{frame}
	
	% ============================================================
	% SLIDE 25: HARD QUESTION 19 - REFERRAL ARRANGEMENTS
	% ============================================================
	\begin{frame}{Hard Question 19: Referral Arrangements}
		\begin{block}{The Scenario (CFA Handbook, p. 204)}
			James Handley receives compensation for each referral he makes to Central Trust's brokerage and personal financial management departments. He does not disclose this to his clients.
		\end{block}
		
		\vspace{0.2cm}
		\textbf{Questions:}
		\begin{enumerate}
			\item Does Handley violate Standard VI(C)?
			\item What if the referral arrangement is \textbf{internal} to the firm?
			\item What if the client would have been referred anyway?
			\item Must disclosure be in writing?
		\end{enumerate}
		
		\vspace{0.2cm}
		\textcolor{cfablue}{\textbf{CFA Answer}}: Violates Standard VI(C). Standard does not distinguish between internal and external payments. Must disclose nature and value of benefit in writing at time of referral.
	\end{frame}
	
	% ============================================================
	% SLIDE 26: HARD QUESTION 20 - WHISTLEBLOWING
	% ============================================================
	\begin{frame}{Hard Question 20: Whistleblowing and Loyalty}
		\begin{block}{The Scenario (CFA Handbook, pp. 134-135)}
			Meredith Rasmussen discovers the hedge fund's reported performance does not reflect a market decline. Her supervisor and compliance officer tell her to "stay away from the issue."
		\end{block}
		
		\vspace{0.2cm}
		\textbf{Questions:}
		\begin{enumerate}
			\item Does whistleblowing violate loyalty to employer?
			\item What steps should Rasmussen take?
			\item When is whistleblowing \textbf{required}?
			\item Does Young's framework apply to whistleblowing?
		\end{enumerate}
		
		\vspace{0.2cm}
		\textcolor{cfablue}{\textbf{CFA Answer}}: Does NOT violate Standard IV(A). Whistleblowing is justified when intent is to protect clients or market integrity, not for personal gain. Must follow firm's reporting procedures first.
	\end{frame}
	
	% ============================================================
	% SLIDE 27: SUMMARY - WHERE YOUNG IS CORRECT
	% ============================================================
	\begin{frame}{Where Young Is Correct}
		\begin{enumerate}
			\item \textbf{Regulatory Lag}: Laws inevitably lag behind innovation. Financial products evolve faster than regulations.
			\item \textbf{Law as Baseline}: Standard I(A) requires following the "more strict" standard - legal is the floor, not the ceiling.
			\item \textbf{Principled Dissociation}: Sometimes ethical obligations require leaving employment.
			\item \textbf{Civil Disobedience}: History shows laws can be unjust (e.g., apartheid, Jim Crow).
		\end{enumerate}
		\vspace{0.2cm}
		\begin{center}
			\fcolorbox{cfablue}{white}{\parbox{0.8\textwidth}{\centering \textbf{Young is valid as an ANALYTICAL correction to oversimplified reasoning}}}
		\end{center}
	\end{frame}
	
	% ============================================================
	% SLIDE 28: SUMMARY - WHERE YOUNG REQUIRES CAUTION
	% ============================================================
	\begin{frame}{Where Young Requires Caution}
		\begin{enumerate}
			\item \textbf{Risk of Rationalization}: Principled resistance can become self-serving justification.
			\item \textbf{The Misconduct Standard}: Standard I(D) applies regardless of legality. Dishonesty and fraud are always unethical.
			\item \textbf{Public vs. Private Defiance}: Principled defiance is \textcolor{cfablue}{public} and accepts consequences. Quiet noncompliance is not the same.
			\item \textbf{Democratic Legitimacy}: Hughes argues it's "presumptively immoral" for firms to unilaterally decide which laws to respect.
		\end{enumerate}
		\vspace{0.2cm}
		\begin{center}
			\fcolorbox{cfablue}{white}{\parbox{0.8\textwidth}{\centering \textbf{Young's framework must be applied with CARE, not convenience}}}
		\end{center}
	\end{frame}
	
	% ============================================================
	% SLIDE 29: FRAMEWORK FOR ETHICAL DECISION-MAKING
	% ============================================================
	\begin{frame}{A Framework for Ethical Decision-Making}
		\small
		\begin{enumerate}
			\item \textbf{Identify Stakeholders} - Who is affected?
			\item \textbf{Evaluate Legal Requirements} - What does the law require?
			\item \textbf{Apply the Code \& Standards} - What do ethical standards require?
			\item \textbf{Consider the "More Strict" Standard} - Which obligation is higher?
			\item \textbf{Assess Impact on Market Integrity} - How does the action affect trust?
			\item \textbf{Document the Analysis} - Could you defend this publicly?
		\end{enumerate}
		
		\vspace{0.2cm}
		\begin{block}{The Three Tests}
			\begin{enumerate}
				\item \textbf{Transparency Test}: Would you defend this action publicly?
				\item \textbf{Social Contract Test}: Does this serve the ultimate benefit of society?
				\item \textbf{Stakeholder Test}: Have you considered impact on all stakeholders?
			\end{enumerate}
		\end{block}
	\end{frame}
	
	% ============================================================
	% SLIDE 30: CONCLUSION
	% ============================================================
	\begin{frame}{Conclusion}
		\begin{block}{Key Takeaway}
			"Although the investment world has become a far more complex place since the first publication of the Standards of Practice Handbook, distinguishing right from wrong remains the paramount principle of the Code and Standards."
			- CFA Institute (2014, p. 5)
		\end{block}
		
		\vspace{0.2cm}
		\textbf{For DBA Scholars:}
		\begin{itemize}
			\item Ethical leadership requires \textcolor{cfablue}{more than legal compliance}
			\item Global business ethics faces challenges laws alone cannot address
			\item Culture of integrity must permeate all levels of operations
			\item Legal standards tell you the \textcolor{cfablue}{floor}, not where to stand
		\end{itemize}
	\end{frame}
	
	% ============================================================
	% SLIDE 31: REFERENCES
	% ============================================================
	\begin{frame}{References}
		\small
		\begin{thebibliography}{9}
			\bibitem{brenkert2019}
			Brenkert, G. G. (2019). Mind the gap! The challenges and limits of (global) business ethics. \textit{Journal of Business Ethics, 155}(4), 917-930.
			
			\bibitem{cfai2014}
			CFA Institute. (2014). \textit{Standards of Practice Handbook} (11th ed.). CFA Institute.
			
			\bibitem{friedman1970}
			Friedman, M. (1970). The social responsibility of business is to increase its profits. \textit{New York Times Magazine}, September 13, 32-33, 122-124.
			
			\bibitem{young2019}
			Young, C. (2019). Putting the law in its place: Business ethics and the assumption that illegal implies unethical. \textit{Journal of Business Ethics, 160}(1), 35-51.
			
			\bibitem{wordpress2013}
			WordPress. (2013, March 15). Ethical theory and its application to contemporary business practice. https://ncys82.wordpress.com/2013/03/15/ethical-theory-and-its-application-to-contemporary-business-practice/
		\end{thebibliography}
	\end{frame}
	
	% ============================================================
	% SLIDE 32: THANK YOU
	% ============================================================
	\begin{frame}{Thank You}
		\begin{center}
			\Huge \textcolor{cfablue}{Questions?}
			
			\vspace{0.5cm}
			\Large \textcolor{cfagold}{"Distinguishing right from wrong remains the paramount principle"}
			
			\vspace{0.5cm}
			\normalsize CFA Institute \textit{Standards of Practice Handbook} (2014, p. 5)
		\end{center}
	\end{frame}
	
	% ============================================================
\end{document}